% Maintenance fixture; elementary calculations, not a research-survey evaluation.
% Copy this file, both *-body.tex files and math-review.sty to one build directory.
\documentclass[12pt,reqno]{amsart}
\usepackage{math-review}
\title[Minimizers and iteration]{Minimizers and iteration for a positive quadratic}
\date{}
\begin{document}
\begin{abstract}
A positive-definite quadratic has an explicit unique minimizer. A fixed-step
gradient iteration converges from every initial point in a precise step-size
range. Completing the square and the spectral decomposition prove these
conclusions and identify the failure at the endpoint. This is an elementary
maintenance fixture, not an independently researched survey.
\end{abstract}
\maketitle
\tableofcontents

\section{Introduction}\label{sec:introduction}
For $m\geq1$, a real symmetric matrix $A\in\mathbb R^{m\times m}$ and
$b\in\mathbb R^m$, consider
\[
Q(x)=\tfrac12 x^{\mathsf T}Ax-b^{\mathsf T}x,\qquad x\in\mathbb R^m.
\]
\begin{definition}
The symmetric matrix $A$ is positive definite when
$v^{\mathsf T}Av>0$ for every nonzero $v\in\mathbb R^m$.
\end{definition}
The two questions are to determine the minimizer and to construct it by
iteration from arbitrary initial data.

\begin{theorem}\label{thm:quadratic}
Let $m\geq1$, $A\in\mathbb R^{m\times m}$ and $b\in\mathbb R^m$.
Assume both of the following:
\begin{itemize}
\item $A=A^{\mathsf T}$;
\item $v^{\mathsf T}Av>0$ for every $v\in\mathbb R^m\setminus\{0\}$.
\end{itemize}
Define
\[
Q(x)=\frac12 x^{\mathsf T}Ax-b^{\mathsf T}x,
\qquad x\in\mathbb R^m.
\]
Then $Q$ has the unique minimizer $x_*=A^{-1}b$ on $\mathbb R^m$.

\end{theorem}
With these hypotheses the minimum value is
\begin{equation}\label{eq:minimum}
Q(x_*)=-\tfrac12 b^{\mathsf T}A^{-1}b.
\end{equation}

\begin{theorem}\label{thm:gradient}
Let $m\geq1$, let $A\in\mathbb R^{m\times m}$ be symmetric positive
definite, and let $b\in\mathbb R^m$. Write $\lambda_{\max}(A)$ for its
largest eigenvalue and set $x_*=A^{-1}b$. Suppose
\[
0<\tau<\frac{2}{\lambda_{\max}(A)}.
\]
For every $x_0\in\mathbb R^m$, define
\[
x_{j+1}=x_j-\tau(Ax_j-b),\qquad j=0,1,\ldots.
\]
Then
\[
\|x_j-x_*\|_2\leq q^j\|x_0-x_*\|_2,\qquad
q=\max_{\lambda\in\operatorname{Spec}(A)}|1-\tau\lambda|<1.
\]
Here $\operatorname{Spec}(A)$ is the set of eigenvalues and $\|\cdot\|_2$
is the Euclidean norm. In particular, $x_j\to x_*$.

\end{theorem}
The step-size restriction is material. For $A=(1)$, $b=0$, $x_0=1$ and
$\tau=2$, the iteration is $x_j=(-1)^j$, so it does not converge.

\section{Completing the square}\label{sec:square}
\typeout{PMS-RECALL-BEFORE=\arabic{theorem}}
\begin{theoremrecall}{thm:quadratic}
Let $m\geq1$, $A\in\mathbb R^{m\times m}$ and $b\in\mathbb R^m$.
Assume both of the following:
\begin{itemize}
\item $A=A^{\mathsf T}$;
\item $v^{\mathsf T}Av>0$ for every $v\in\mathbb R^m\setminus\{0\}$.
\end{itemize}
Define
\[
Q(x)=\frac12 x^{\mathsf T}Ax-b^{\mathsf T}x,
\qquad x\in\mathbb R^m.
\]
Then $Q$ has the unique minimizer $x_*=A^{-1}b$ on $\mathbb R^m$.

\end{theoremrecall}
\typeout{PMS-RECALL-AFTER=\arabic{theorem}}
\begin{proof}[Proof of Theorem~\ref{thm:quadratic}]
If $Av=0$, then $v^{\mathsf T}Av=0$. Positive definiteness gives $v=0$,
so $A$ is invertible. Set $x_*=A^{-1}b$. Symmetry gives
\[
Q(x)-Q(x_*)=\tfrac12(x-x_*)^{\mathsf T}A(x-x_*).
\]
The right-hand side is strictly positive for $x\ne x_*$. Substitution of
$x_*$ proves~\eqref{eq:minimum}.
\end{proof}

\section{The spectral contraction}\label{sec:contraction}
\typeout{PMS-GRADIENT-RECALL-BEFORE=\arabic{theorem}}
\begin{theoremrecall}{thm:gradient}
Let $m\geq1$, let $A\in\mathbb R^{m\times m}$ be symmetric positive
definite, and let $b\in\mathbb R^m$. Write $\lambda_{\max}(A)$ for its
largest eigenvalue and set $x_*=A^{-1}b$. Suppose
\[
0<\tau<\frac{2}{\lambda_{\max}(A)}.
\]
For every $x_0\in\mathbb R^m$, define
\[
x_{j+1}=x_j-\tau(Ax_j-b),\qquad j=0,1,\ldots.
\]
Then
\[
\|x_j-x_*\|_2\leq q^j\|x_0-x_*\|_2,\qquad
q=\max_{\lambda\in\operatorname{Spec}(A)}|1-\tau\lambda|<1.
\]
Here $\operatorname{Spec}(A)$ is the set of eigenvalues and $\|\cdot\|_2$
is the Euclidean norm. In particular, $x_j\to x_*$.

\end{theoremrecall}
\typeout{PMS-GRADIENT-RECALL-AFTER=\arabic{theorem}}
\begin{proof}[Proof of Theorem~\ref{thm:gradient}]
Let $e_j=x_j-x_*$. Since $Ax_*=b$,
\[
e_{j+1}=(I-\tau A)e_j,\qquad e_j=(I-\tau A)^j e_0.
\]
The spectral theorem gives an orthogonal matrix $U$ with
$U^{\mathsf T}AU=\operatorname{diag}(\lambda_1,\ldots,\lambda_m)$.
Each $\lambda_i>0$, and the step-size condition implies
$-1<1-\tau\lambda_i<1$. Consequently
\[
\|I-\tau A\|_2=\max_i|1-\tau\lambda_i|=q<1.
\]
Taking norms in the error formula proves the estimate and convergence.
\end{proof}
For the one-dimensional data in the introduction, the error multiplier is
$1-\tau=-1$. This proves the stated endpoint failure; it does not claim
that every datum must fail outside the guaranteed range.

\section{A checked example}\label{sec:example}
Let
\[
A=\begin{pmatrix}2&0\\0&4\end{pmatrix},\qquad
b=\begin{pmatrix}2\\8\end{pmatrix},\qquad \tau=\tfrac14.
\]
The hypotheses are directly verifiable:
\begin{enumerate}
\item $A=A^{\mathsf T}$ and $v^{\mathsf T}Av=2v_1^2+4v_2^2>0$ for $v\ne0$.
\item $0<\tau=1/4<2/\lambda_{\max}(A)=1/2$.
\end{enumerate}
Thus
\[
x_*=\begin{pmatrix}1\\2\end{pmatrix},\qquad Q(x_*)=-9,\qquad
q=\max\{|1-2\tau|,|1-4\tau|\}=\tfrac12.
\]
For every $x_0$, Theorem~\ref{thm:gradient} gives
$\|x_j-x_*\|_2\leq2^{-j}\|x_0-x_*\|_2$.
\end{document}
